\section{The exact tree deformation scheme}\label{sec:flat}

Let \(G\) be a finite connected tree with at least two vertices. At vertex \(v\), take \(S^1\) times a sphere with \(\deg(v)+2\) boundary circles. Denote its regular fibre by \(h_v\), and use the zero-framing section \(z\) at each torus. At the two surplus tori fill slopes \(-2z+h\) and \(2z+h\). Along each graph edge glue by the section--fibre exchange, with both signs \(\epsilon_{vw}=\epsilon_{wv}\in\{1,-1\}\) in the marked torus coordinates. The edge maps reverse the boundary orientation. This defines a closed oriented 3-manifold \(M_G\), with the same convention as~\cite{Paper48,Paper49}. In the fundamental group we label the two cone generators so that
\[
        h_v=U_v^{-2}=V_v^{2}.
\]

Fix a flat representation with every \(h_v=-1\) and nonparallel axes at the endpoints of every seam. In a spanning-tree gauge (here the graph itself), all transported quaternions are in one common copy of \(\SU\). The local angular coordinates use a fixed oriented frame modulo ordinary conjugacy. We regard \(\mathfrak{su}(2)\) as imaginary quaternions and write dot products in its unit quaternion basis.

\begin{theorem}[Exact Jacobian scheme]\label{thm:flat}
Near the reference conjugacy class the local analytic flat deformation scheme is the critical scheme of
\begin{equation}\label{eq:treephi}
     \Phi_G(p,T)=\sum_{vw\in E(G)}T_{vw}p_vp_w.
\end{equation}
Here \(p_v\) is defined by \(U_v=u_v\exp(p_vu_v)\), \(u_v\in S^2\), and \(T_{vw}\) are analytic edge coordinates equal at \(p=0\) to the signed dot \(B_{vw}=\epsilon_{vw}u_v\cdot u_w\). There are also passive angular coordinates. The assertion is an equality of analytic ideals in the indicated local chart, retaining nilpotent deformations.
\end{theorem}
\begin{proof}
Choose the vertex boundary loops so that the section product has the two cone factors first. At \(p=0\) the other section factors are central. After extracting their fixed central signs, their ordered product \(g_v\) is
\[
        g_v=\prod_{w\sim v}^{\text{boundary order}}\exp(-2\epsilon_{vw}p_wu_w).
\]
The vertex product eliminates \(V_v\) analytically as a fixed central sign times \(U_v^{-1}g_v^{-1}\). Since the relation \(V_v^2=U_v^{-2}\) is unaffected by that sign, it is equivalent to \(U_vg_vU_v^{-1}=g_v^{-1}\). All other commutations with \(V_v\) follow from those with \(U_v\) and the individual section factors. The commutations with the individual section factors are \([h_v,h_w]=1\), since sections are neighboring fibres up to the seam sign, and the commutation with \(U_v\) follows from \(h_v=U_v^{-2}\).

The difference \(h_vh_w-h_wh_v\) is \(2\sin(2p_v)\sin(2p_w)(u_v\times u_w)\). One component of the cross product is a unit at a nonparallel reference point. Hence the commuting ideal is
\[
                  J_0=(p_vp_w:vw\in E(G)).
\]
The analytic logarithm \(L_v=\log g_v\) has each component in the ideal generated by \(p_w\) with \(w\sim v\). The remaining inversion equation is \((\Ad_{U_v}+I)L_v=0\). Modulo \(J_0\), \(p_v L_v=0\), so \(\Ad_{U_v}\) can be replaced by \(\Ad_{u_v}\), the rotation by \(\pi\) about \(u_v\). It follows that the only remaining equation at a vertex is the scalar \(F_v=-u_v\cdot L_v/2=0\).

Analytic radial integration in the neighbor variables gives
\[
 F_v=\sum_{w\sim v}p_w(B_{vw}+c_{vw}),\qquad
              c_{vw}\in(p_j:j\sim v).
\]
Set \(T_{vw}=B_{vw}+c_{vw}+c_{wv}\). The omitted contribution \(p_wc_{wv}\) is in \(J_0\), so the full ideal is exactly generated by \(p_vp_w\) and \(\sum_w T_{vw}p_w\), the partial derivatives of \eqref{eq:treephi}. Boundary-order commutator terms are included in the analytic \(c_{vw}\).

To see that the \(B_{vw}\) are independent local coordinates modulo passive angles, fix the root axis and its first neighbor in an oriented frame by conjugation. Each further child axis is determined by its dot with its parent and one local azimuth; the first edge already fixes the frame. The angular derivative of \(T\) equals that of \(B\) at \(p=0\). Thus \(T\) is an analytic coordinate change. Every elimination above holds in the analytic local ring, and hence also over finite real local Artin algebras.
\end{proof}

\begin{lemma}[Comparison with the analytic Coulomb reduction]\label{lem:coulomb}
Let \(g\) be the potential \eqref{eq:treephi} after analytically splitting off any nondegenerate Hessian directions at the chosen point. The actual stationary reduction \(f\) of Chern--Simons to harmonic variables can be expressed in analytic coordinates in which \(\Jac_f=\Jac_g\), where \(\Jac_h\) denotes the ideal generated by all partial derivatives of \(h\).
\end{lemma}
\begin{proof}
Nonparallel neighboring axes give \(H^0(M_G;\operatorname{ad}A)=0\) at the reference flat connection \(A\); its stabilizer is the finite center, which acts trivially. Use a smooth metric and Sobolev order \(s>5/2\). The fixed Coulomb slice decomposes into \(a\in\mathcal H^1_A\) and its coexact complement \(b\). On the complement \(*d_A\) is invertible from order \(s\) to \(s-1\), as follows from the Hodge Laplacian. The analytic implicit function theorem solves the coexact stationary equation by \(b=b(a)=O(a^2)\). Set \(f(a)=\CS(A+a+b(a))-\CS(A)\) for a local real branch. Its stationary equation is the remaining harmonic projection of \(*F_{A+a+b(a)}=0\). When this holds, the only possible component is \(*F_{A+a+b(a)}=d_A\phi\). Bianchi gives \(d_{A+a+b(a)}^*d_A\phi=0\), and \(\phi=0\) by the small perturbation of the invertible scalar Laplacian. The same reasoning applies formally over each finite Artin extension of the reference point, where inversion on nilpotent perturbations is finite.

For an actual analytic comparison map, take the holonomies of \(A+a+b(a)\) around the selected cone loops \emph{off} the flat locus, use the same angular frame and the coordinates of Theorem~\ref{thm:flat}, and project onto the variables remaining after the graph-model Hessian splitting. Path restriction at this regularity, the linear path ODE and the coordinate operations are analytic. On Artin critical points this map is the flat-bundle/holonomy equivalence. In this use of the holonomy construction, one forms the flat bundle on the universal cover from its representation over the Artin algebra and identifies the underlying smooth lift by the nilpotent filtration and partitions of unity. Conversely, nilpotent flat parallel transport recovers the representation; smooth gauge corresponds to basepoint conjugation. The Coulomb gauge is unique in each infinitesimal neighborhood by inversion of \(d_A^*d_A\), order by order in the nilpotent ideal. The construction of the underlying groupoid equivalence is the general holonomy step of~\cite[Sec.~6]{GoldmanMillson}; no K\"ahler formality theorem is invoked here.

On dual numbers the analytic map is the period isomorphism from \(H^1\) to infinitesimal holonomy modulo conjugacy. After the indicated Hessian splitting the graph-model Jacobian has no linear part, so the retained coordinates have the full tangent dimension. The analytic map is a local diffeomorphism. In those same coordinates, the equivalence over all finite local Artin rings identifies the completed Jacobian ideals (the functors are those represented by the completed local quotients). Their contractions to convergent germs are equal by the Krull intersection theorem applied to the analytic local quotients. This proves the stated ideal identity, rather than equality of radicals only.
\end{proof}
